\chapter{Lumbar Drain Placement}

\section*{Overview}
A lumbar drain is a small catheter inserted into the lumbar subarachnoid space to continuously drain cerebrospinal fluid (CSF) to lower intracranial pressure or manage CSF leaks.

\section*{Alternative Names}
Spinal catheter placement, Continuous CSF drainage

\section*{Principle}
A catheter is inserted through a needle between the lumbar vertebrae into the subarachnoid space and connected to a closed drainage system. It allows controlled, continuous drainage of cerebrospinal fluid to lower intracranial pressure or remove blood and debris.
\section*{History}
Evolved from the lumbar puncture technique first described by Heinrich Quincke in 1891. Continuous lumbar drainage was developed in the mid-20th century.

\section*{Why It Is Done}
Ordered to treat postoperative CSF leaks, reduce brain tension during neurosurgery, or evaluate patients for normal pressure hydrocephalus (NPH).

\section*{Is This a Primary or Secondary Test or Procedure}
Primary temporary therapeutic option for persistent cranial/spinal CSF leaks.

\section*{How It Is Performed}
Under sterile conditions, a needle is inserted into the L3-L4 or L4-L5 space. A flexible catheter is threaded through the needle into the subarachnoid space and connected to a collection bag.

\section*{Preparation}
Check coagulation profile (coagulopathy is a strict contraindication). Position the patient sitting or in lateral decubitus.

\section*{Contraindications and Precautions}
Obstructive hydrocephalus or space-occupying lesion with mass effect (risk of brain herniation), or local skin infection.

\section*{Risks}
Over-drainage (headache, subdural hematoma, herniation), meningitis, nerve root irritation (radicular pain), or catheter blockage.

\section*{Aftercare and Recovery}
The drainage system must be leveled carefully (usually at the shoulder or ear level) as ordered. Monitor neurological status and drainage rate hourly.

\section*{Outcome and Follow-up}
Clear CSF is collected. Symptoms of CSF leak should resolve.

\section*{Expected Results}
Drainage of clear, colorless CSF at the prescribed flow rate (typically 5-10 mL/hour); patient neurologically intact.

\section*{Follow-up}
Daily CSF laboratory analysis (cell count, protein, glucose) if infection is suspected; catheter removal within 5-7 days.

\section*{When to Seek Medical Care}
Follow the procedure team's specific instructions. Seek urgent medical assessment for severe or worsening pain, uncontrolled bleeding, fever or signs of infection, trouble breathing, chest pain, fainting, new weakness or confusion, inability to urinate, or any rapidly worsening symptom. The appropriate warning signs depend on the procedure and the patient's medical condition.

\section*{References}
\begin{enumerate}
  \item MedlinePlus. Lumbar Drain Placement. U.S. National Library of Medicine; 2025.
  \item Current Medical Diagnosis and Treatment. McGraw-Hill; 2025.
\end{enumerate}

\clearpage
